\chapter{CAD建模技术}
\label{ch:cad-modeling}
\label{ch:modeling}

\section{从外形表达走向设计意图}
CAD不仅是把手工图纸搬到计算机上。二维图纸表达视图与尺寸，三维模型直接组织空间形状、装配关系，并为CAE和CAM提供数据基础。讲义概括的技术演进为：\textbf{线框造型$\to$曲面造型$\to$实体造型$\to$参数化$\to$变量化}。

线框模型用顶点和棱边描述骨架，数据较少，但缺乏面的信息，可能产生形状二义性；曲面模型增加曲面及其边界与连接关系，适合自由外形；实体模型进一步明确边界哪一侧为实体内部，使点的内外判断、体积与质量等计算有了基础。

参数化方法用尺寸和约束驱动模型修改，并通过特征间的父子关系传播变化；变量化方法侧重把几何关系组织为约束方程并求解，不必完全依赖严格的顺序构造。实际系统可以结合多种方法，不能简单用软件名称划出互斥类别。
\sourcepage{第3章}{3--14、37--39}

\section{几何、拓扑与有效实体}
\subsection{形体由哪些层次组成}
\dfn{几何与拓扑}{\textbf{几何信息}给出点的坐标、曲线与曲面方程，以及位置、尺寸和方向；\textbf{拓扑信息}给出元素的连接、邻接、包含与方向关系。例如“边位于某条圆曲线上”是几何关联，“该边由两个顶点界定并属于两个相邻面”是拓扑关联。}

顶点对应端点位置；边是被端点界定的曲线段；有序、有向边构成封闭的\textbf{环}；面是在支撑曲面上由外环及若干内环界定的区域；相邻面形成\textbf{壳}，外壳界定外边界，内壳可界定空腔；体则表示被封闭边界围住的三维区域。讲义中的\textbf{基本体/构造体素}是由有限尺寸参数定义的简单实体，如长方体、球和圆柱，应与离散网格中的微小体积单元区分。
\sourcepage{第3章}{16--17}

\subsection{正则化集合运算}
用于实体建模的对象应有界、维数均匀、边界内外可判定，不能把悬点、悬边或悬面混入三维实体。课程采用的闭合流形边界还要求边的邻接关系有效，普通边在闭合二流形边界上邻接两个面。

设$\operatorname{int}$为取内部、$\operatorname{cl}$为取闭包，正则化为
\[
r(X)=\operatorname{cl}(\operatorname{int}X).
\]
正则并、交、差定义为
\[
A\cup^{*}B=r(A\cup B),\quad
A\cap^{*}B=r(A\cap B),\quad
A-^{*}B=r(A-B).
\]
先去掉无三维内部的降维成分，再恢复三维区域的边界。例如两个实体仅以一个面接触，普通交集是该面，正则交集为空，而不是把零厚度面当作实体。
\sourcepage{第3章}{18--19}

\subsection{Euler关系：先识别形体及计数对象}
对单个连通、无贯通孔、无内部空腔，且各面为单连通区域的闭合多面体，
\[
V-E+F=2.
\]
$V,E,F$分别为顶点、边、面数。长方体为$8-12+6=2$。若面含内环，记全部面的内环总数为$L$；对无内部空腔的$B$个独立实体、总贯通孔数$G$，课程采用
\[
V-E+F-L=2(B-G).
\]
若有独立内壳，应按各闭合壳分别计算Euler特征，再相加，不能继续把$B$简单取为实体个数。内环数与贯通孔数对应不同层次：盲孔底面、通孔和封闭空腔也不能混淆。

讲义另给$V-E+F-B=1$，用于把一个连通、无洞的区域划分为$B$个三维胞腔的计数，此时$F$\textbf{包含内部公共面}，$B$是胞腔数。它与$B$表示互不相连的独立实体时的公式不能混用。

\nt{Euler关系是拓扑计数的必要检查，不能单独证明模型有效。即使计数成立，仍可能有自交、面朝向错误、未闭合边界或几何不一致。}
\sourcepage{第3章}{20--22}

\section{常见建模操作}
\begin{center}
\renewcommand{\arraystretch}{1.3}
\begin{tabularx}{\linewidth}{lX}
\toprule
操作 & 构造思路及结果 \\
\midrule
基本体与布尔运算 & 创建球、圆柱、长方体等，再作并、交、差，组合出复杂实体。 \\
拉伸、旋转、扫描 & 使截面沿直线、绕轴或沿路径运动；闭合区域在适当条件下可生成实体，开放轮廓通常生成曲面。 \\
放样/蒙皮 & 让一组有序截面之间平滑过渡；若未封闭端面，结果可能仅为曲面。 \\
混合/圆角 & 用光滑过渡面替代尖锐边或顶点，改善连接处形状。 \\
抬升与扭曲 & 局部移动或替换面、边、顶点，改变已有边界。 \\
边界建模 & 直接增删、修改顶点、边和面；修改几何时仍需维护有效拓扑。 \\
特征建模 & 用孔、槽、凹腔、圆角、倒角等有工程含义的单元构建模型，并保留附加信息。 \\
\bottomrule
\end{tabularx}
\end{center}
\sourcepage{第3章}{23--35}

\subsection{特征与参数为什么比单纯外形更有用}
一个圆柱形空腔既可以仅作为被减去的几何体，也可以作为“孔特征”，带有孔径、深度、定位、螺纹等语义。后者更容易关联加工方法和检验要求。特征模型不仅保存最终形状，还组织设计意图、工艺和管理信息。

参数化草图的一般过程是：画出大致轮廓，施加几何约束与尺寸，求解并重建二维图形，再由拉伸、旋转等生成三维形体，必要时建立尺寸间的关系。\textbf{欠约束}表示尚有可自由变化的量；\textbf{过约束}可能是约束重复或冲突；修改参数前要了解依赖关系。

\ex{孔距如何体现设计意图？}{若对称布置的两个孔应随板宽变化始终居中，可定义板中线为基准，以对称约束和孔距控制位置。若仅把两个孔的位置分别固定为绝对坐标，修改板宽后未必保持对称。几何形状暂时相同，设计关系可能不同。}
\sourcepage{第3章}{32--36}

\section{实体的几种数据表达}
\subsection{B-Rep：边界表达法}
\dfn{B-Rep（Boundary Representation）}{通过顶点、边、环、面、壳及其连接关系表达实体边界，并为各拓扑元素关联几何数据。面关联曲面，边关联曲线，顶点关联坐标；边界方向用来区分实体内外。}

\reviewfig[.9]{ch03-brep.pdf}{边界模型中几何信息与拓扑信息的对应（摘自第3章讲义，PDF第46页）。}{fig:brep}

曲面方程本身可能定义无限延伸的支撑曲面，面则只占其上被环裁定的有限区域；曲线与边也具有同样区别。只保存曲线曲面方程不足以知道它们如何构成封闭实体，只保存连接表又不足以确定实际位置和形状。

顶点表、边表和面表是基本组织方式；链表允许元素数量动态变化，循环双向链表便于沿环正反遍历。有孔的面应增加内环，有内部空腔的实体应增加内壳，不能用不存在的几何桥接边来替代真实边界。
\sourcepage{第3章}{38--48}

\subsection{半边结构：同一条边在两个面中分别使用}
一条共享边在两个相邻面的边界序列中，其“下一条边”通常不同，不能给整条边只存一个统一的后继。半边结构把一条边拆为\textbf{方向相反的两个半边}，每个半边服务于一个面的一个环。

\begin{center}
\renewcommand{\arraystretch}{1.25}
\begin{tabularx}{\linewidth}{lX}
\toprule
元素 & 主要关联 \\
\midrule
体 & 所属面、边、顶点的列表。 \\
面 & 所属体、外环、内环列表、曲面几何。 \\
环 & 所属面、任一半边。 \\
边 & 两个配对的半边。 \\
半边 & 所属边、起始顶点、所属环、前驱与后继。 \\
顶点 & 坐标、从该顶点出发的一个半边。 \\
\bottomrule
\end{tabularx}
\end{center}

\reviewfig[.48]{ch03-loops.pdf}{面的外环与内环（摘自第3章讲义，PDF第51页）。}{fig:halfedge-loops}

沿面选定法向观察时，外环与内环取相反方向，使沿有向边行进时面区域始终在左侧。环的方向用于保证面朝向和内外区域判定一致。

\subsubsection*{常见邻接查询怎样做}
\begin{enumerate}
\item \textbf{一个面的边与顶点：}从外环的任一半边出发，按\texttt{next}遍历到起点，收集其所属边与起始顶点；再对每个内环重复。
\item \textbf{一条边的相邻面：}分别取两个半边，经“半边$\to$环$\to$面”获得。
\item \textbf{一个顶点的关联边与面：}在闭合流形网格中，从出发半边$h$开始，重复“先取配对半边，再取其\texttt{next}”，绕该顶点一周，逐个收集所属边与面。
\end{enumerate}
有开放边界时需处理缺失的配对半边或显式边界半边；非流形顶点也可能需要多个环。查询规则成立的前提是数据结构满足相应的拓扑假设。
\sourcepage{第3章}{49--56}

\subsection{翼边结构与Euler操作}
翼边以一条有向边为中心，记录起终顶点、左右相邻面（或环），以及左右边界上相邻的四条“翼边”。它与半边都通过额外邻接信息加快局部查询，区别在于组织的基本单元及指针安排。

\reviewfig[.87]{ch03-winged-edge.pdf}{翼边所保存的局部邻接关系（摘自第3章讲义，PDF第57页）。}{fig:winged-edge}

修改B-Rep时，几何与拓扑都可能变化。Euler操作是建立或修改顶点、边、面等连接关系的基本拓扑操作，设计时使元素数量的变化符合Euler关系；实际建模还需保证曲线曲面的几何一致性及实体有效性。
\sourcepage{第3章}{57--59}

\subsection{CSG：保存构造过程的树}
\dfn{CSG（Constructive Solid Geometry）}{由简单基本体及正则布尔运算构造实体。CSG树的叶节点保存基本体及其参数、位姿，内部节点保存并、交或差运算，树结构保存组合顺序。}

\reviewfig[.9]{ch03-csg.pdf}{实体$(P_0-P_1)\cup P_2$及其CSG构造树（摘自第3章讲义，PDF第60页）。}{fig:csg}

CSG表达紧凑，便于修改基本体参数，也便于保存构造逻辑；不足是局部边界修改不直观，能够方便构造的形状受基本体和运算限制，显示前还需计算边界。B-Rep直接保存边界，适合显示与局部几何处理；系统常同时保存特征/构造历史和B-Rep结果，两者分工而非互相排斥。
\sourcepage{第3章}{60--61}

\subsection{体素、八叉树与单元分解}
体素模型将空间划分为小立方体，用占据状态近似表示物体。若边长为$h$，且各占据单元质量均匀，体积可估计为$V\approx N_{\rm occ}h^3$。分辨率越高，表面逼近越细，但均匀网格每个方向的单元数翻倍，总数就约增至8倍。

八叉树对空间立方体递归分为8个子块：全空、全满的块可直接标记；同时含内外区域的混合块继续细分。因此可对边界附近保留细节，并减少均匀区域的存储。单元分解进一步用不重叠的简单单元组合复杂形体；它强调空间的分块，而B-Rep强调边界，CSG强调布尔构造。
\sourcepage{第3章}{62--64}

\section{装配模型与约束}
\dfn{装配模型}{同时包含产品组成结构、各零部件几何与设计参数、装配层次，以及连接、配合、约束和运动关系的产品模型。装配树描述“由什么组成”，装配约束描述“如何放置及如何连接”。}

常见约束包括贴合、对齐、距离、角度、相切等。应由实际功能选择约束并检查自由度：两零件外观上摆在一起，不表示已建立稳定的装配关系；约束也不一定消除所有相对运动，例如转动副有意保留一个转动自由度。
\sourcepage{第3章}{66--68}

\begin{center}
\renewcommand{\arraystretch}{1.3}
\begin{tabularx}{\linewidth}{lXX}
\toprule
比较 & 自底向上 & 自顶向下 \\
\midrule
顺序 & 先完成独立零件，再定义装配关系 & 先确定总体结构、基准与关系，再细化零部件 \\
关联 & 零件设计相对独立，适合现有件与标准件 & 零件可引用总体或相邻件，尺寸和结构关联更强 \\
优点 & 思路直观，独立设计和复用方便 & 便于总体协调及并行设计，减少不匹配造成的返工 \\
代价 & 缺乏前期统筹时容易不满足装配要求 & 引用关系需管理，系统和设计人员要求较高 \\
\bottomrule
\end{tabularx}
\end{center}

例如先按总体轮距和安装接口建立万向轮装配，再在环境中设计支架和轮体，属于自顶向下；从零件库取现有轮体与支架并用配合约束组装，属于自底向上。实际项目常混合使用。
\sourcepage{第3章}{69--77}

\section{MBD：把制造信息组织到三维模型中}
\dfn{MBD（Model Based Definition）}{基于模型的定义以三维模型组织和定义产品信息，将传统二维图纸中的尺寸、公差、粗糙度、技术要求等产品制造信息PMI（Product and Manufacturing Information）与几何模型关联，使模型成为产品定义与下游应用的数据载体。}

MBD数据包括三部分：\textbf{设计模型}，如几何形状、基准、辅助几何和尺寸；\textbf{注解}，如尺寸公差、粗糙度、焊接符号、明细表和技术要求；\textbf{属性}，即不直接显示但补充模型定义的产品信息。

信息可以通过特征树、层、注解视图及其组合管理。讲义示例在注解文件夹中组织不同视图，借助标注工具表达尺寸公差，再按3D PDF、eDrawings或STEP 242等形式发布。复习重点是\textbf{几何与语义信息的完整传递}；具体接收系统是否能读取、保留可计算语义，还取决于导出设置和软件支持，不能只凭扩展名判断。

\sourcepage{第3章}{79--83}

\Needspace{9cm}
\subsection*{复习速查与自测}
\begin{enumerate}
\item 顶点坐标与“该边属于哪两个面”分别是什么信息？\textit{答：几何信息与拓扑信息。}
\item 曲面方程与一个面是否相同？\textit{答：不同，面还需要有限的裁定区域、边界及方向。}
\item 为什么半边比只给整条边存一个后继更合理？\textit{答：同一边在不同面中的后继不同，半边分别记录各面的边界次序。}
\item CSG树的叶节点和内部节点分别存什么？\textit{答：基本体及参数位姿、布尔运算。}
\item 体素边长减半，均匀三维网格存储量大约如何变化？\textit{答：约增至8倍。}
\item 单个无孔、无空腔长方体的Euler关系是什么？\textit{答：$8-12+6=2$。}
\item MBD比只含外形的三维模型增加了什么？\textit{答：与几何关联的尺寸、公差、工艺注解及附加属性等产品定义信息。}
\end{enumerate}
